% Gemeinsame Präambel für Skript.tex und die Einzelkapitel (einzelkapitel.tex)
% Gemeinsamer Stil der Übungsblätter und Tests (Kopie im Ordner stil/)
\makeatletter\def\input@path{{stil/}}\makeatother
\input{gdi-stil}
\input{header}
\usepackage{needspace}
\usepackage{multicol}
\usetikzlibrary{shapes.geometric, shapes.multipart, positioning}

%%%%%%%%%%%%%%%%%%%%%%%%%%%%%%%%%%%%%%%%%%%%%%%%%%
% Bausteine des Skripts
%%%%%%%%%%%%%%%%%%%%%%%%%%%%%%%%%%%%%%%%%%%%%%%%%%
% Merkkasten (wie "Kurz erklärt" auf den Arbeitsblättern)
\newtcolorbox{merke}[1][Merke]{enhanced, breakable, colback=diagramm!5, colframe=diagramm!80, boxrule=0.6pt, arc=1.5mm,
	left=3mm, right=3mm, top=1.5mm, bottom=1.5mm, title=#1, fonttitle=\bfseries\small, colbacktitle=diagramm!80, coltitle=white,
	before skip=3mm, after skip=3mm}
% Hinweis-Kasten in den Farben der Code-Kästen
\newtcolorbox{hinweis}{enhanced, sharp corners, boxrule=0pt, colback=codeHintergrund,
	borderline west={2pt}{0pt}{diagramm}, left=3mm, top=1.5mm, bottom=1.5mm, before skip=3mm, after skip=3mm}

% Die drei Teile jedes Kapitels
\newcommand{\teil}[1]{\needspace{9\baselineskip}\subsection*{\textcolor{codeSchluessel}{#1}}}
\newcommand{\worum}{\teil{Worum geht es?}}
\newcommand{\beispiele}{\teil{Beispiele}}
\newcommand{\aufgaben}{\teil{Aufgaben}}
\newcommand{\ueben}{\teil{Zum Weiterüben}}

% Beispiel mit Titel
\newcounter{beispiel}[section]
\newcommand{\beispiel}[1]{\par\medskip\needspace{6\baselineskip}\refstepcounter{beispiel}%
	\noindent\textbf{\textcolor{diagramm}{Beispiel \thesection.\arabic{beispiel}:} #1}\par\nopagebreak\smallskip}

% Aufgabe mit Titel; optionale Markierung: \K = wie in der Klausur, \S = anspruchsvoll
\newcounter{aufgabe}[section]
\renewcommand{\theaufgabe}{\thesection.\arabic{aufgabe}}
\newcommand{\K}{\ \tcbox[on line, sharp corners, boxrule=0pt, colback=diagramm!80, coltext=white, left=1mm, right=1mm, top=0.2mm, bottom=0.2mm, fontupper=\sffamily\bfseries\scriptsize]{KLAUSUR}}
\renewcommand{\S}{\ \textcolor{diagramm}{$\bigstar$}}
\newcommand{\aufgabe}[2][]{\par\medskip\needspace{4\baselineskip}\refstepcounter{aufgabe}%
	\noindent\textbf{Aufgabe \theaufgabe: #2}#1\par\nopagebreak\smallskip}

% Tabellen
\newcolumntype{C}[1]{>{\centering\arraybackslash}p{#1}}
\newcolumntype{L}[1]{>{\raggedright\arraybackslash}p{#1}}
\newcommand{\leer}{\rule[-0.6em]{0pt}{2.0em}}
\newcommand{\luecke}[1][3cm]{\rule[-0.4ex]{#1}{0.4pt}}
\newcommand{\code}[1]{\texttt{#1}}
% In Schreibmaschinenschrift keine Ligaturen (-- bleibt --) und " bleibt ein Anführungszeichen
\DisableLigatures{encoding = *, family = tt*}
\AtBeginDocument{\shorthandoff{"}}

% Flussdiagramme (wie in den Tests zu Blatt 4)
\tikzset{
	startstop/.style={rectangle, rounded corners=3mm, draw=diagramm, semithick, fill=diagramm!15, minimum width=2cm, minimum height=7mm, font=\small},
	io/.style={trapezium, trapezium left angle=75, trapezium right angle=105, draw=diagramm, semithick, fill=white, minimum height=7mm, font=\small, align=center},
	process/.style={rectangle, draw=diagramm, semithick, fill=white, minimum width=2.6cm, minimum height=7mm, font=\small, align=center},
	decision/.style={diamond, aspect=2.4, draw=diagramm, semithick, fill=diagramm!8, inner sep=1pt, font=\small, align=center},
	arrow/.style={-{Latex[length=2mm]}, draw=diagramm, thick},
	lbl/.style={font=\footnotesize, fill=white, inner sep=1.5pt},
	% UML-Klassen
	klasse/.style={rectangle split, rectangle split parts=3, draw=diagramm, semithick, fill=diagramm!5, font=\ttfamily\footnotesize, align=left, rectangle split part align={center,left,left}, minimum width=3.2cm},
	erbt/.style={-{Triangle[open, length=3mm, width=3mm]}, draw=diagramm, semithick},
	assoz/.style={draw=diagramm, semithick}
}

% Kopf- und Fußzeile
\clearpairofpagestyles
\ihead{GDI -- Kompaktskript}
\ohead{\rightmark}
\cfoot{\pagemark}
\automark{section}
\setkomafont{pageheadfoot}{\normalfont\small\color{codeKommentar}}
\setcounter{tocdepth}{1}
\emergencystretch=3em
